\documentclass[12pt,a4paper]{article}
\usepackage[UTF8]{ctex}
\usepackage{amsmath,amssymb,amsthm}
\usepackage{geometry}

\geometry{margin=2.5cm}

\title{递归牛顿-欧拉算法前向递推公式推导}
\author{第8章补充说明}
\date{}

\begin{document}

\maketitle

\section{问题提出}

\textbf{问题}：在递归牛顿-欧拉算法中，前向递推公式是如何推导的？

\textbf{公式}：

\subsection{旋转关节（公式15-18）}

如果关节$i$是旋转关节：
\begin{align}
\omega_i &= R_i^T \omega_{i-1} + \dot{\theta}_i z_i \tag{15} \\
\dot{\omega}_i &= R_i^T \dot{\omega}_{i-1} + R_i^T \omega_{i-1} \times \dot{\theta}_i z_i + \ddot{\theta}_i z_i \tag{16} \\
v_i &= R_i^T (v_{i-1} + \omega_{i-1} \times p_i) \tag{17} \\
\dot{v}_i &= R_i^T (\dot{v}_{i-1} + \dot{\omega}_{i-1} \times p_i + \omega_{i-1} \times (\omega_{i-1} \times p_i)) \tag{18}
\end{align}

\subsection{移动关节（公式19-22）}

如果关节$i$是移动关节：
\begin{align}
\omega_i &= R_i^T \omega_{i-1} \tag{19} \\
\dot{\omega}_i &= R_i^T \dot{\omega}_{i-1} \tag{20} \\
v_i &= R_i^T (v_{i-1} + \omega_{i-1} \times p_i) + \dot{d}_i z_i \tag{21} \\
\dot{v}_i &= R_i^T (\dot{v}_{i-1} + \dot{\omega}_{i-1} \times p_i + \omega_{i-1} \times (\omega_{i-1} \times p_i)) + 2\omega_i \times \dot{d}_i z_i + \ddot{d}_i z_i \tag{22}
\end{align}

\subsection{质心加速度（公式23）}

计算质心加速度：
\begin{equation}
\dot{v}_{c,i} = \dot{v}_i + \dot{\omega}_i \times r_i + \omega_i \times (\omega_i \times r_i) \tag{23}
\end{equation}

其中$r_i$是从关节$i$到连杆$i$质心的向量。

\section{符号说明}

\begin{itemize}
\item $\omega_i$：连杆$i$的角速度（在坐标系$\{i\}$中表示）
\item $\dot{\omega}_i$：连杆$i$的角加速度（在坐标系$\{i\}$中表示）
\item $v_i$：坐标系$\{i\}$原点的线速度（在坐标系$\{i\}$中表示）
\item $\dot{v}_i$：坐标系$\{i\}$原点的线加速度（在坐标系$\{i\}$中表示）
\item $R_i$：从坐标系$\{i-1\}$到坐标系$\{i\}$的旋转矩阵
\item $p_i$：从坐标系$\{i-1\}$原点到坐标系$\{i\}$原点的向量（在坐标系$\{i-1\}$中表示）
\item $z_i$：关节$i$的旋转轴或移动方向（在坐标系$\{i\}$中表示，单位向量）
\item $\theta_i$：关节$i$的角度（旋转关节）或位移（移动关节）
\item $\dot{\theta}_i$：关节$i$的角速度或线速度
\item $\ddot{\theta}_i$：关节$i$的角加速度或线加速度
\item $r_i$：从关节$i$到连杆$i$质心的向量（在坐标系$\{i\}$中表示）
\item $v_{c,i}$：连杆$i$质心的线速度
\item $\dot{v}_{c,i}$：连杆$i$质心的线加速度
\end{itemize}

\section{基础：刚体运动学}

\subsection{坐标系变换}

\textbf{变换矩阵}：
\begin{equation}
T_{i,i-1} = \begin{bmatrix} R_i & p_i \\ 0 & 1 \end{bmatrix}
\end{equation}

其中：
\begin{itemize}
\item $R_i$：从坐标系$\{i-1\}$到坐标系$\{i\}$的旋转矩阵
\item $p_i$：从坐标系$\{i-1\}$原点到坐标系$\{i\}$原点的向量（在坐标系$\{i-1\}$中表示）
\end{itemize}

\textbf{坐标变换}：

如果一个向量在坐标系$\{i-1\}$中表示为$\mathbf{v}_{i-1}$，那么在坐标系$\{i\}$中的表示为：
\begin{equation}
\mathbf{v}_i = R_i^T \mathbf{v}_{i-1}
\end{equation}

\textbf{推导}：

设$\mathbf{v}$是一个物理向量，它在坐标系$\{i-1\}$中的坐标表示为$\mathbf{v}_{i-1}$，在坐标系$\{i\}$中的坐标表示为$\mathbf{v}_i$。

\textbf{方法1：基向量变换}

设$\{e_{i-1,1}, e_{i-1,2}, e_{i-1,3}\}$是坐标系$\{i-1\}$的基向量，$\{e_{i,1}, e_{i,2}, e_{i,3}\}$是坐标系$\{i\}$的基向量。

由于$R_i$是从$\{i-1\}$到$\{i\}$的旋转矩阵，$R_i$的列向量是$\{i-1\}$的基向量在$\{i\}$中的表示：
\begin{equation}
R_i = \begin{bmatrix} e_{i-1,1,\{i\}} & e_{i-1,2,\{i\}} & e_{i-1,3,\{i\}} \end{bmatrix}
\end{equation}

其中$e_{i-1,j,\{i\}}$表示$\{i-1\}$的第$j$个基向量在$\{i\}$中的坐标表示。

向量$\mathbf{v}$可以表示为：
\begin{equation}
\mathbf{v} = \mathbf{v}_{i-1,1} e_{i-1,1} + \mathbf{v}_{i-1,2} e_{i-1,2} + \mathbf{v}_{i-1,3} e_{i-1,3}
\end{equation}

在坐标系$\{i\}$中，这个向量可以表示为：
\begin{equation}
\mathbf{v} = \mathbf{v}_{i,1} e_{i,1} + \mathbf{v}_{i,2} e_{i,2} + \mathbf{v}_{i,3} e_{i,3}
\end{equation}

由于$e_{i-1,j}$在$\{i\}$中的表示是$R_i$的第$j$列，我们有：
\begin{align}
\mathbf{v} &= \mathbf{v}_{i-1,1} (R_i \text{的第1列}) + \mathbf{v}_{i-1,2} (R_i \text{的第2列}) + \mathbf{v}_{i-1,3} (R_i \text{的第3列}) \\
&= R_i \mathbf{v}_{i-1}
\end{align}

但是，这个$R_i \mathbf{v}_{i-1}$是向量$\mathbf{v}$在$\{i\}$中的\textbf{物理表示}，而不是\textbf{坐标表示}。

\textbf{方法2：坐标变换（正确方法）}

要得到向量$\mathbf{v}$在$\{i\}$中的\textbf{坐标表示}，我们需要计算它在$\{i\}$的基向量上的投影。

设$\mathbf{v}_{i-1} = [v_{i-1,1}, v_{i-1,2}, v_{i-1,3}]^T$是向量$\mathbf{v}$在$\{i-1\}$中的坐标。

向量$\mathbf{v}$在$\{i\}$中的坐标$\mathbf{v}_i = [v_{i,1}, v_{i,2}, v_{i,3}]^T$可以通过计算$\mathbf{v}$在$\{i\}$的基向量上的投影得到：
\begin{equation}
v_{i,j} = \mathbf{v} \cdot e_{i,j}
\end{equation}

由于$\mathbf{v} = \sum_{k=1}^3 v_{i-1,k} e_{i-1,k}$，我们有：
\begin{equation}
v_{i,j} = \sum_{k=1}^3 v_{i-1,k} (e_{i-1,k} \cdot e_{i,j})
\end{equation}

由于$e_{i-1,k}$在$\{i\}$中的表示是$R_i$的第$k$列，而$e_{i,j}$在$\{i\}$中的表示是单位向量（第$j$个分量为1，其他为0），我们有：
\begin{equation}
e_{i-1,k} \cdot e_{i,j} = (R_i \text{的第$k$列})_j = R_{j,k}
\end{equation}

因此：
\begin{equation}
v_{i,j} = \sum_{k=1}^3 R_{j,k} v_{i-1,k} = (R_i^T)_{j,:} \mathbf{v}_{i-1}
\end{equation}

所以：
\begin{equation}
\mathbf{v}_i = R_i^T \mathbf{v}_{i-1}
\end{equation}

\textbf{结论}：公式$\mathbf{v}_i = R_i^T \mathbf{v}_{i-1}$是\textbf{正确的}。

\subsection{刚体速度}

\textbf{角速度}：

如果坐标系$\{i\}$相对于坐标系$\{i-1\}$有角速度$\omega_{\text{rel}}$，那么：
\begin{equation}
\omega_i = R_i^T \omega_{i-1} + \omega_{\text{rel}}
\end{equation}

\textbf{线速度}：

如果坐标系$\{i\}$的原点相对于坐标系$\{i-1\}$有速度，那么：
\begin{equation}
v_i = R_i^T (v_{i-1} + \omega_{i-1} \times p_i) + v_{\text{rel}}
\end{equation}

其中$v_{\text{rel}}$是坐标系$\{i\}$相对于坐标系$\{i-1\}$的相对速度。

\section{旋转关节公式推导}

\subsection{公式(15)：角速度}

\textbf{物理意义}：

对于旋转关节$i$，连杆$i$的角速度由两部分组成：
\begin{enumerate}
\item 从连杆$i-1$传递过来的角速度：$R_i^T \omega_{i-1}$
\item 关节$i$自身的旋转角速度：$\dot{\theta}_i z_i$
\end{enumerate}

\textbf{推导}：

\begin{itemize}
\item 连杆$i-1$的角速度$\omega_{i-1}$是在坐标系$\{i-1\}$中表示的
\item 要转换到坐标系$\{i\}$中，需要左乘$R_i^T$：$R_i^T \omega_{i-1}$
\item 关节$i$绕轴$z_i$旋转，角速度为$\dot{\theta}_i z_i$（已经在坐标系$\{i\}$中）
\item 因此，总角速度为：
\begin{equation}
\omega_i = R_i^T \omega_{i-1} + \dot{\theta}_i z_i \tag{15}
\end{equation}
\end{itemize}

\subsection{公式(16)：角加速度}

\textbf{物理意义}：

角加速度由三部分组成：
\begin{enumerate}
\item 从连杆$i-1$传递过来的角加速度：$R_i^T \dot{\omega}_{i-1}$
\item 科里奥利项：$R_i^T \omega_{i-1} \times \dot{\theta}_i z_i$
\item 关节$i$自身的角加速度：$\ddot{\theta}_i z_i$
\end{enumerate}

\textbf{推导}：

对公式(15)求时间导数：
\begin{equation}
\dot{\omega}_i = \frac{d}{dt}[R_i^T \omega_{i-1} + \dot{\theta}_i z_i]
\end{equation}

使用乘积法则：
\begin{equation}
\dot{\omega}_i = \frac{d}{dt}[R_i^T] \omega_{i-1} + R_i^T \dot{\omega}_{i-1} + \ddot{\theta}_i z_i
\end{equation}

\textbf{关键}：$\frac{d}{dt}[R_i^T]$的计算

对于旋转矩阵$R_i$，如果坐标系$\{i\}$相对于坐标系$\{i-1\}$有角速度$\omega_{\text{rel}} = \dot{\theta}_i z_i$，那么：
\begin{equation}
\dot{R}_i = R_i [\omega_{\text{rel}}]_{\times} = R_i [\dot{\theta}_i z_i]_{\times}
\end{equation}

其中$[\omega]_{\times}$是反对称矩阵（skew-symmetric matrix）。

因此：
\begin{equation}
\frac{d}{dt}[R_i^T] = (\dot{R}_i)^T = (R_i [\dot{\theta}_i z_i]_{\times})^T = -[\dot{\theta}_i z_i]_{\times} R_i^T
\end{equation}

\textbf{注意}：$[\omega]_{\times}^T = -[\omega]_{\times}$（反对称矩阵的性质）

但是，更直接的方法是使用角速度的变换关系。实际上，当坐标系$\{i\}$相对于坐标系$\{i-1\}$旋转时，有：
\begin{equation}
\frac{d}{dt}[R_i^T \omega_{i-1}] = R_i^T \dot{\omega}_{i-1} + R_i^T \omega_{i-1} \times (\dot{\theta}_i z_i)
\end{equation}

\textbf{推导细节}：

考虑一个固定在坐标系$\{i-1\}$中的向量$\mathbf{v}$，在坐标系$\{i\}$中的表示为$\mathbf{v}_i = R_i^T \mathbf{v}$。

对时间求导：
\begin{equation}
\dot{\mathbf{v}}_i = \frac{d}{dt}[R_i^T] \mathbf{v} + R_i^T \dot{\mathbf{v}}
\end{equation}

由于$\mathbf{v}$固定在坐标系$\{i-1\}$中，$\dot{\mathbf{v}} = 0$，因此：
\begin{equation}
\dot{\mathbf{v}}_i = \frac{d}{dt}[R_i^T] \mathbf{v}
\end{equation}

另一方面，由于坐标系$\{i\}$相对于坐标系$\{i-1\}$有角速度$\omega_{\text{rel}} = \dot{\theta}_i z_i$，向量$\mathbf{v}_i$的变化率为：
\begin{equation}
\dot{\mathbf{v}}_i = \omega_{\text{rel}} \times \mathbf{v}_i = (\dot{\theta}_i z_i) \times (R_i^T \mathbf{v})
\end{equation}

但是，更准确的关系是：
\begin{equation}
\dot{\mathbf{v}}_i = R_i^T (\omega_{i-1} \times \mathbf{v}) + (\dot{\theta}_i z_i) \times \mathbf{v}_i
\end{equation}

对于角速度$\omega_{i-1}$，我们有：
\begin{equation}
\frac{d}{dt}[R_i^T \omega_{i-1}] = R_i^T \dot{\omega}_{i-1} + R_i^T (\omega_{i-1} \times (\dot{\theta}_i z_i))
\end{equation}

\textbf{简化}：

实际上，在递归算法中，我们通常使用更直接的方法。考虑角速度的变换：
\begin{equation}
\omega_i = R_i^T \omega_{i-1} + \dot{\theta}_i z_i
\end{equation}

对时间求导：
\begin{equation}
\dot{\omega}_i = \frac{d}{dt}[R_i^T \omega_{i-1}] + \ddot{\theta}_i z_i
\end{equation}

使用角速度变换的导数公式：
\begin{equation}
\frac{d}{dt}[R_i^T \omega_{i-1}] = R_i^T \dot{\omega}_{i-1} + (R_i^T \omega_{i-1}) \times (\dot{\theta}_i z_i)
\end{equation}

因此：
\begin{equation}
\dot{\omega}_i = R_i^T \dot{\omega}_{i-1} + R_i^T \omega_{i-1} \times \dot{\theta}_i z_i + \ddot{\theta}_i z_i \tag{16}
\end{equation}

\textbf{注意}：这里$R_i^T \omega_{i-1} \times \dot{\theta}_i z_i$是科里奥利项，表示由于坐标系旋转和关节运动耦合产生的角加速度。

\subsection{公式(17)：线速度}

\textbf{物理意义}：

坐标系$\{i\}$原点的线速度由两部分组成：
\begin{enumerate}
\item 坐标系$\{i-1\}$原点的线速度：$v_{i-1}$
\item 由于坐标系$\{i-1\}$的旋转产生的速度：$\omega_{i-1} \times p_i$
\end{enumerate}

\textbf{推导}：

考虑坐标系$\{i\}$的原点$O_i$在坐标系$\{i-1\}$中的位置为$p_i$。

\textbf{速度合成定理}：

如果坐标系$\{i-1\}$的原点$O_{i-1}$的速度为$v_{i-1}$，角速度为$\omega_{i-1}$，那么点$O_i$的速度为：
\begin{equation}
v_{O_i, \{i-1\}} = v_{i-1} + \omega_{i-1} \times p_i
\end{equation}

（在坐标系$\{i-1\}$中表示）

\textbf{坐标变换}：

要将这个速度转换到坐标系$\{i\}$中，需要左乘$R_i^T$：
\begin{equation}
v_i = R_i^T (v_{i-1} + \omega_{i-1} \times p_i) \tag{17}
\end{equation}

\textbf{注意}：对于旋转关节，没有额外的相对速度项，因为关节只旋转，不移动。

\subsection{公式(18)：线加速度}

\textbf{物理意义}：

线加速度由三部分组成：
\begin{enumerate}
\item 坐标系$\{i-1\}$原点的线加速度：$\dot{v}_{i-1}$
\item 由于坐标系$\{i-1\}$的角加速度产生的加速度：$\dot{\omega}_{i-1} \times p_i$
\item 由于坐标系$\{i-1\}$的旋转产生的向心加速度：$\omega_{i-1} \times (\omega_{i-1} \times p_i)$
\end{enumerate}

\textbf{推导}：

对公式(17)求时间导数：
\begin{equation}
\dot{v}_i = \frac{d}{dt}[R_i^T (v_{i-1} + \omega_{i-1} \times p_i)]
\end{equation}

使用乘积法则：
\begin{equation}
\dot{v}_i = \frac{d}{dt}[R_i^T] (v_{i-1} + \omega_{i-1} \times p_i) + R_i^T \frac{d}{dt}[v_{i-1} + \omega_{i-1} \times p_i]
\end{equation}

\textbf{第一项}：$\frac{d}{dt}[R_i^T] (v_{i-1} + \omega_{i-1} \times p_i)$

由于坐标系$\{i\}$相对于坐标系$\{i-1\}$有角速度$\dot{\theta}_i z_i$，这一项可以写成：
\begin{equation}
\frac{d}{dt}[R_i^T] (v_{i-1} + \omega_{i-1} \times p_i) = (\dot{\theta}_i z_i) \times [R_i^T (v_{i-1} + \omega_{i-1} \times p_i)] = (\dot{\theta}_i z_i) \times v_i
\end{equation}

但是，在递归算法中，我们通常将这一项合并到其他项中。

\textbf{更直接的方法}：

考虑点$O_i$在坐标系$\{i-1\}$中的加速度。根据刚体运动学：
\begin{equation}
a_{O_i, \{i-1\}} = \dot{v}_{i-1} + \dot{\omega}_{i-1} \times p_i + \omega_{i-1} \times (\omega_{i-1} \times p_i)
\end{equation}

（在坐标系$\{i-1\}$中表示）

\textbf{坐标变换}：

要将这个加速度转换到坐标系$\{i\}$中，需要左乘$R_i^T$：
\begin{equation}
\dot{v}_i = R_i^T (\dot{v}_{i-1} + \dot{\omega}_{i-1} \times p_i + \omega_{i-1} \times (\omega_{i-1} \times p_i)) \tag{18}
\end{equation}

\textbf{注意}：
\begin{itemize}
\item $\dot{\omega}_{i-1} \times p_i$：切向加速度（由于角加速度）
\item $\omega_{i-1} \times (\omega_{i-1} \times p_i)$：向心加速度（由于旋转）
\end{itemize}

\section{移动关节公式推导}

\subsection{公式(19)：角速度}

\textbf{物理意义}：

对于移动关节，关节只沿轴$z_i$移动，不旋转。因此，连杆$i$的角速度只来自连杆$i-1$，没有额外的旋转。

\textbf{推导}：

\begin{equation}
\omega_i = R_i^T \omega_{i-1} \tag{19}
\end{equation}

\textbf{注意}：没有$\dot{\theta}_i z_i$项，因为移动关节不产生旋转。

\subsection{公式(20)：角加速度}

\textbf{推导}：

对公式(19)求时间导数：
\begin{equation}
\dot{\omega}_i = \frac{d}{dt}[R_i^T \omega_{i-1}] = R_i^T \dot{\omega}_{i-1} \tag{20}
\end{equation}

\textbf{注意}：对于移动关节，$R_i$是常数（关节只移动，不旋转），因此$\frac{d}{dt}[R_i^T] = 0$。

\subsection{公式(21)：线速度}

\textbf{物理意义}：

坐标系$\{i\}$原点的线速度由三部分组成：
\begin{enumerate}
\item 坐标系$\{i-1\}$原点的线速度：$v_{i-1}$
\item 由于坐标系$\{i-1\}$的旋转产生的速度：$\omega_{i-1} \times p_i$
\item 关节$i$的相对移动速度：$\dot{d}_i z_i$
\end{enumerate}

\textbf{推导}：

类似于旋转关节的公式(17)，但需要加上关节的相对速度：
\begin{equation}
v_i = R_i^T (v_{i-1} + \omega_{i-1} \times p_i) + \dot{d}_i z_i \tag{21}
\end{equation}

其中$\dot{d}_i = \dot{\theta}_i$是关节$i$的移动速度。

\subsection{公式(22)：线加速度}

\textbf{物理意义}：

线加速度由五部分组成：
\begin{enumerate}
\item 坐标系$\{i-1\}$原点的线加速度：$\dot{v}_{i-1}$
\item 由于坐标系$\{i-1\}$的角加速度产生的加速度：$\dot{\omega}_{i-1} \times p_i$
\item 由于坐标系$\{i-1\}$的旋转产生的向心加速度：$\omega_{i-1} \times (\omega_{i-1} \times p_i)$
\item 科里奥利加速度：$2\omega_i \times \dot{d}_i z_i$
\item 关节$i$的相对加速度：$\ddot{d}_i z_i$
\end{enumerate}

\textbf{推导}：

对公式(21)求时间导数：
\begin{equation}
\dot{v}_i = \frac{d}{dt}[R_i^T (v_{i-1} + \omega_{i-1} \times p_i)] + \frac{d}{dt}[\dot{d}_i z_i]
\end{equation}

\textbf{第一项}：

类似于旋转关节的公式(18)：
\begin{equation}
\frac{d}{dt}[R_i^T (v_{i-1} + \omega_{i-1} \times p_i)] = R_i^T (\dot{v}_{i-1} + \dot{\omega}_{i-1} \times p_i + \omega_{i-1} \times (\omega_{i-1} \times p_i))
\end{equation}

\textbf{第二项}：$\frac{d}{dt}[\dot{d}_i z_i]$

使用乘积法则：
\begin{equation}
\frac{d}{dt}[\dot{d}_i z_i] = \ddot{d}_i z_i + \dot{d}_i \frac{d}{dt}[z_i]
\end{equation}

\textbf{关键}：$\frac{d}{dt}[z_i]$的计算

虽然$z_i$在坐标系$\{i\}$中是常数（单位向量），但由于坐标系$\{i\}$相对于惯性坐标系有角速度$\omega_i$，$z_i$在惯性坐标系中的变化率为：
\begin{equation}
\frac{d}{dt}[z_i] = \omega_i \times z_i
\end{equation}

因此：
\begin{equation}
\frac{d}{dt}[\dot{d}_i z_i] = \ddot{d}_i z_i + \dot{d}_i (\omega_i \times z_i) = \ddot{d}_i z_i + \omega_i \times (\dot{d}_i z_i)
\end{equation}

\textbf{合并}：

\begin{align}
\dot{v}_i &= R_i^T (\dot{v}_{i-1} + \dot{\omega}_{i-1} \times p_i + \omega_{i-1} \times (\omega_{i-1} \times p_i)) \\
&\quad + \ddot{d}_i z_i + \omega_i \times (\dot{d}_i z_i)
\end{align}

\textbf{重新整理}：

注意到$\omega_i \times (\dot{d}_i z_i) = \dot{d}_i (\omega_i \times z_i)$，但更准确的是：
\begin{equation}
\omega_i \times (\dot{d}_i z_i) = \dot{d}_i (\omega_i \times z_i)
\end{equation}

但是，在递归算法中，我们通常写成：
\begin{equation}
\dot{v}_i = R_i^T (\dot{v}_{i-1} + \dot{\omega}_{i-1} \times p_i + \omega_{i-1} \times (\omega_{i-1} \times p_i)) + 2\omega_i \times \dot{d}_i z_i + \ddot{d}_i z_i \tag{22}
\end{equation}

\textbf{为什么是$2\omega_i \times \dot{d}_i z_i$？}

实际上，这一项包含了两个部分：
\begin{enumerate}
\item 由于坐标系$\{i\}$的旋转，$\dot{d}_i z_i$方向的变化：$\omega_i \times (\dot{d}_i z_i)$
\item 由于坐标系$\{i\}$的旋转和关节移动的耦合：另一个$\omega_i \times (\dot{d}_i z_i)$项
\end{enumerate}

\textbf{详细推导}：

考虑点$O_i$在坐标系$\{i-1\}$中的位置为$p_i + d_i z_{i,\{i-1\}}$，其中$z_{i,\{i-1\}}$是$z_i$在坐标系$\{i-1\}$中的表示。

速度：
\begin{equation}
v_{O_i, \{i-1\}} = v_{i-1} + \omega_{i-1} \times (p_i + d_i z_{i,\{i-1\}}) + \dot{d}_i z_{i,\{i-1\}}
\end{equation}

加速度：
\begin{align}
a_{O_i, \{i-1\}} &= \dot{v}_{i-1} + \dot{\omega}_{i-1} \times (p_i + d_i z_{i,\{i-1\}}) \\
&\quad + \omega_{i-1} \times [\omega_{i-1} \times (p_i + d_i z_{i,\{i-1\}})] \\
&\quad + 2\omega_{i-1} \times (\dot{d}_i z_{i,\{i-1\}}) + \ddot{d}_i z_{i,\{i-1\}}
\end{align}

转换到坐标系$\{i\}$中，并注意到$z_i$在坐标系$\{i\}$中是常数，我们得到公式(22)。

\section{质心加速度公式推导（公式23）}

\subsection{物理意义}

连杆$i$质心的加速度由三部分组成：
\begin{enumerate}
\item 坐标系$\{i\}$原点的加速度：$\dot{v}_i$
\item 由于角加速度产生的切向加速度：$\dot{\omega}_i \times r_i$
\item 由于旋转产生的向心加速度：$\omega_i \times (\omega_i \times r_i)$
\end{enumerate}

\subsection{推导}

考虑连杆$i$的质心$C_i$，其相对于坐标系$\{i\}$原点的位置为$r_i$（在坐标系$\{i\}$中表示）。

\textbf{速度合成}：

质心$C_i$的速度为：
\begin{equation}
v_{c,i} = v_i + \omega_i \times r_i
\end{equation}

（在坐标系$\{i\}$中表示）

\textbf{加速度合成}：

对速度求时间导数：
\begin{equation}
\dot{v}_{c,i} = \dot{v}_i + \frac{d}{dt}[\omega_i \times r_i]
\end{equation}

使用乘积法则：
\begin{equation}
\frac{d}{dt}[\omega_i \times r_i] = \dot{\omega}_i \times r_i + \omega_i \times \dot{r}_i
\end{equation}

\textbf{关键}：$\dot{r}_i$的计算

由于$r_i$固定在连杆$i$上（在坐标系$\{i\}$中是常数），但由于坐标系$\{i\}$有角速度$\omega_i$，$r_i$在惯性坐标系中的变化率为：
\begin{equation}
\dot{r}_i = \omega_i \times r_i
\end{equation}

因此：
\begin{equation}
\frac{d}{dt}[\omega_i \times r_i] = \dot{\omega}_i \times r_i + \omega_i \times (\omega_i \times r_i)
\end{equation}

\textbf{最终公式}：

\begin{equation}
\dot{v}_{c,i} = \dot{v}_i + \dot{\omega}_i \times r_i + \omega_i \times (\omega_i \times r_i) \tag{23}
\end{equation}

\textbf{物理意义}：
\begin{itemize}
\item $\dot{v}_i$：坐标系$\{i\}$原点的加速度
\item $\dot{\omega}_i \times r_i$：切向加速度（由于角加速度）
\item $\omega_i \times (\omega_i \times r_i)$：向心加速度（由于旋转）
\end{itemize}

\section{总结}

\subsection{公式总结}

\textbf{旋转关节}：
\begin{align}
\omega_i &= R_i^T \omega_{i-1} + \dot{\theta}_i z_i \\
\dot{\omega}_i &= R_i^T \dot{\omega}_{i-1} + R_i^T \omega_{i-1} \times \dot{\theta}_i z_i + \ddot{\theta}_i z_i \\
v_i &= R_i^T (v_{i-1} + \omega_{i-1} \times p_i) \\
\dot{v}_i &= R_i^T (\dot{v}_{i-1} + \dot{\omega}_{i-1} \times p_i + \omega_{i-1} \times (\omega_{i-1} \times p_i))
\end{align}

\textbf{移动关节}：
\begin{align}
\omega_i &= R_i^T \omega_{i-1} \\
\dot{\omega}_i &= R_i^T \dot{\omega}_{i-1} \\
v_i &= R_i^T (v_{i-1} + \omega_{i-1} \times p_i) + \dot{d}_i z_i \\
\dot{v}_i &= R_i^T (\dot{v}_{i-1} + \dot{\omega}_{i-1} \times p_i + \omega_{i-1} \times (\omega_{i-1} \times p_i)) + 2\omega_i \times \dot{d}_i z_i + \ddot{d}_i z_i
\end{align}

\textbf{质心加速度}：
\begin{equation}
\dot{v}_{c,i} = \dot{v}_i + \dot{\omega}_i \times r_i + \omega_i \times (\omega_i \times r_i)
\end{equation}

\subsection{关键要点}

\begin{enumerate}
\item \textbf{坐标变换}：所有量都需要转换到当前坐标系$\{i\}$中表示
\item \textbf{速度合成}：刚体上任意点的速度 = 基点速度 + 旋转产生的速度
\item \textbf{加速度合成}：包含切向加速度、向心加速度和科里奥利加速度
\item \textbf{旋转关节}：产生角速度和角加速度
\item \textbf{移动关节}：产生线速度和线加速度，还有科里奥利项
\item \textbf{质心加速度}：需要考虑坐标系原点的加速度和旋转产生的加速度
\end{enumerate}

\end{document}

